\section{Related Works}

Modern LLM training typically begins with pretraining on massive datasets, with the goal of modeling the natural language distribution and acquiring broad world knowledge. Mid-training generally follows the same underlying learning paradigm, while differences in training data and parameter configurations can lead to substantially different outcomes. In the post-training stage, a common practice is to train specialized expert models for tasks in different domains, typically using SFT for cold start followed by task-specific RL algorithms. These expert models can then be distilled into a unified model through methods such as OPD, resulting in a final general-purpose model~\citep{deepseek-ai_deepseek-v4_2026, team_kimi_2026, glm-5-team_glm-5_2026}.

\subsection{Distribution Matching Perspective}

Among the algorithms used across different stages of LLM training, most admit a relatively clear interpretation from the perspective of distribution matching, with RL being a notable exception. Despite its distinctive optimization formulation, RL has also been widely recognized as closely related to distribution matching~\citep{korbak_reinforcement_2022}. Prior to our work, a number of studies investigated the connection between RL and distribution-matching objectives and derived locally optimal solutions with an exponential-tilting form similar to that in our Theorem~\ref{lem:reward-reweighting}~\citep{korbak_rl_2022, korbak_reinforcement_2022}. Reward-induced conditional distributions have likewise been studied in prior work and have motivated several algorithmic improvements~\citep{tajwar_maximum_2026, zhou_variational_2025}.

In most previous studies, however, such distributions primarily appeared as auxiliary results in the derivation or analysis of RL objectives. Subtle discrepancies remain between RL optimization and direct distribution matching, preventing the two perspectives from being fully unified~\citep{korbak_reinforcement_2022}. Recent studies on MaxRL provide an important step toward bridging this gap, revealing a first-order approximation between RL and distribution matching. More importantly, our work differs in that we treat the target distribution induced by distribution matching as a central object of study, rather than merely an intermediate analytical construct. In particular, we study both its role in estimating the training objective and its stochastic evolution throughout training, which together constitute the main results of this work. We expect that this perspective may provide a useful lens for further understanding and improving related training algorithms.

\subsection{Entropy Collapse in RL}

Diversity collapse and entropy decrement in reinforcement learning have long been recognized in the classical RL literature. RL optimization can drive a model toward a suboptimal solution with reduced behavioral diversity, thereby limiting the model's long-term training potential. Classical approaches to mitigating this issue include entropy regularization~\citep{ahmed_understanding_2019} and constraints based on KL divergence~\citep{schulman_trust_2017}. These methods, however, introduce their own trade-offs, including optimization instability and restrictions on reward improvement.

As RL has become increasingly prevalent in LLM post-training, a growing body of work has revisited entropy collapse in this setting. Existing studies have investigated its causes and potential remedies from a variety of perspectives, including biases introduced by PPO clipping~\citep{yu_dapo_2025, cai_beyond_2026, jin_revisiting_2026}, geometric biases induced by model architecture~\citep{ren_learning_2025}, and the underlying training dynamics~\citep{wang_entropy_2026}. In contrast, our analysis identifies a long-horizon degradation of the target distribution caused by variance accumulation. This phenomenon arises at a more fundamental level and does not depend on a particular RL algorithm or implementation trick. Based on this analysis, we further derive a effective regularization strategy. Entropy collapse may represent one of the central challenges in scaling post-training, and our results provide an additional perspective on its underlying mechanism.

\subsection{Variance Control and Training Stability}

The stability of RL algorithms has been studied extensively, with the high variance of training-objective estimators being one of the primary challenges. A substantial fraction of algorithmic developments in RL can be understood as attempts to address this issue~\citep{konda_onactor-critic_2003, schulman_trust_2017, schulman_proximal_2017}. Advantage estimation is among the most widely adopted approaches: it can be implemented using a separate critic model~\citep{konda_onactor-critic_2003, schulman_trust_2017, schulman_proximal_2017}, or approximated through techniques such as group-based sampling and normalization~\citep{shao_deepseekmath_2024, ahmadian_back_2024}. Beyond advantage estimation, the treatment of importance weights and the design of the underlying training infrastructure can also have a substantial impact on optimization stability~\citep{zheng_group_2025, gao_soft_2025}.

Recent empirical evidence suggests that, in the LLM setting, the strong prior induced by extensive pretraining can enable stable RL training over a meaningful range even in the absence of aggressive variance-reduction techniques~\citep{ahmadian_back_2024}. Nevertheless, controlling training variance can still yield measurable performance gains. Our analysis further provides a possible explanation for these gains: reducing training variance not only stabilizes individual optimization steps, but can also mitigate the degradation of the training target caused by variance accumulation over time. As a result, the model is more likely to preserve meaningful exploration and reach better regions of the parameter space.
